\begin{exercice}
Le chiffre des unités pour l'écriture décimale  $3^{100}$ est~:
\begin{enumerate}
  \item $3$
  \item $9$
  \item $1$
  \item $7$
  \item $6$
\end{enumerate}
\end{exercice}

\begin{proof} %[\textbf{C}]
Rappelons que la division euclienne d'un entier positif $N$ par $10$ fournit les entiers $q$ (le quotient) et $r$ (le reste) %
tels que: 
%
\begin{equation}
  N = q \cdot 10  + r \qquad (0 \leq r < 10).
\end{equation}
%
Rechercher le chiffre des unités revient donc à trouver le reste $r$. Pour ce faire, nous partons de 
%
\begin{equation}
  3^{100} = (3^2)^{50} = 9^{50} =(9^2)^{25}.
\end{equation}
%
Or 
\begin{equation}
  9^2 = 81 \equiv 1 \mod 10.
\end{equation}
%
D'où
%
\begin{equation}
  3^{100} = (9^2)^{25} \equiv 1^{25} \equiv 1 \mod 10.
\end{equation}
%
Alternativement, on peut aussi noter que 
%
\begin{equation}
  3^{100} = 9^{50} \equiv (\minus 1)^{50} \equiv 1 \mod 10. 
\end{equation}
%
Réponse \textbf{C}.
%
\end{proof}